\section{小数据与双语学习：把人类直觉变成可检验实验}

课程借用儿童学习提出数据效率问题，但并未把神经网络等同于大脑。人类获得视觉、行动、社会互动、长期目标和高结构输入；语言模型通常获得离线 token 与统一 loss。研究价值在于构造受控数据条件，而不是用少量 token 直接证明“类脑学习”。

\subsection{为什么研究 small model 与 small data}

Baby Scale 对比 infant/ChatGPT 与 brain/neural network，列出连续学习、grounded input、目标、层级结构和潜在 inductive bias。这里的对比是 hypothesis generator：如果儿童输入中交互性、语义多样性或说话人结构与网页不同，可以设计 dataset-level test；但不能把人类全部监督折算成转录文本字符。

\lecturefigure{slide-039.jpg}{儿童与 ChatGPT、大脑与神经网络：两类学习系统的对照}{CS25 V6 official deck, p.~39}

\lecturefigure{slide-040.jpg}{人类学习与 LLM 训练在连续性、目标、模态和结构上的差异}{CS25 V6 official deck, p.~40}

\lecturefigure{slide-041.jpg}{研究小模型与小数据的动机：效率、解释性与受控实验}{CS25 V6 official deck, p.~41}

\teachervoice{00:18:22--00:20:46，老师反复使用 might、potential 与 question 来描述人类学习优势。课堂提示：这些差异是实验轴，不是“儿童数据天然优于互联网”的结论。}

\begin{warningbox}{比较儿童与 LLM token 数时缺失了什么}
儿童同时接收图像、声音、动作后果、社会反馈和反复互动；转录语料只保留其中一小部分。模型与儿童的目标函数、参数初始化、先验和评估任务也不同。因此 token 数可用于提出数据效率问题，却不是公平的总监督量。
\end{warningbox}

\subsection{Baby Scale：一个孩子的语言输入能训练出什么}

Baby Scale 使用 CHILDES 中单个家庭的 child-directed speech，训练小型 GPT-2 风格模型，并在 language modeling、syntax、semantics 与 lexical development 等任务上比较。核心设计不是追求 SOTA，而是固定模型规模后观察\textbf{家庭数据来源}与\textbf{儿童接触结构}是否带来可测差异。

\lecturefigure{slide-043.jpg}{Baby Scale 核心问题：模型能否从一个孩子的输入中学习}{CS25 V6 official deck, p.~43}

\lecturefigure{slide-044.jpg}{实验设置：家庭转录、token 规模、模型与评估任务}{CS25 V6 official deck, p.~44}

Perplexity（困惑度）是交叉熵的指数：
\[
\operatorname{PPL}=\exp\!\left(-\frac{1}{T}\sum_{t=1}^{T}\log p_{\theta}(x_t\mid x_{<t})\right).
\]
直觉上它近似模型在每个 token 面前的“有效候选数”；越低表示对测试序列越不惊讶。但 PPL 受 tokenizer、测试分布和长度影响，不能跨任意模型直接比较，也不等同于语法、事实或交互能力。

结果页应按两步读：slide 45 比较不同家庭数据在多项任务上的相对表现，slide 46 再看模型大小是否改变趋势。若某个数据集在多项任务都较强，可能来自词汇覆盖、互动结构或说话人多样性；如果扩大模型不能消除差异，说明 data property 不是单纯 capacity bottleneck。

\lecturefigure{slide-045.jpg}{Baby Scale 多任务结果：不同家庭输入产生不同能力轮廓}{CS25 V6 official deck, p.~45}

\lecturefigure{slide-046.jpg}{模型规模与数据来源的交互：扩大参数并不统一消除差异}{CS25 V6 official deck, p.~46}

\begin{knowledgebox}{读图：先比较任务，再比较规模}
先在每个 subplot 内比较不同家庭数据，而不是跨任务比较绝对分数；各任务 metric 和难度不同。再看模型大小变化后排序是否稳定。稳定排序支持“数据属性持续影响表现”，却不能仅凭相关性确定是 child interaction、词汇覆盖还是转录质量导致，需要后续受控 ablation。
\end{knowledgebox}

\subsection{Scale 之外：哪些数据属性预测表现}

课程列出 interaction、speaker diversity、context length、vocabulary coverage、child/caregiver speech ratio 与 distribution divergence 等候选变量。可写成一个简单回归视角：
\[
y_{f,m}=\beta_0+\beta_1\log D_f+\beta_2\log N_m
+\sum_k\gamma_k z_{f,k}+\varepsilon_{f,m},
\]
$y_{f,m}$ 是家庭 $f$、模型规模 $m$ 的任务表现，$D_f$ 是 token 数，$N_m$ 是参数量，$z_{f,k}$ 是数据属性。该式不是论文的唯一模型，而是帮助读者区分 size effect 与 composition effect。

\lecturefigure{slide-047.jpg}{Scaling results：数据集与模型规模对不同任务的影响并不一致}{CS25 V6 official deck, p.~47}

\lecturefigure{slide-048.jpg}{Beyond size：互动、词汇、说话人和上下文等候选解释变量}{CS25 V6 official deck, p.~48}

\lecturefigure{slide-049.jpg}{总体结论：表现较好的家庭数据更丰富、更具互动与语义结构}{CS25 V6 official deck, p.~49}

\lecturefigure{slide-050.jpg}{Baby Scale takeaways：child-scale learning 可行但高度 task-dependent}{CS25 V6 official deck, p.~50}

\teachervoice{00:22:23--00:25:25，老师强调“更大”不是唯一解释：child/caregiver 比例、互动多样性与词汇结构可能影响表现；同时不同任务给出的最优数据集并不完全一致。}

\begin{importantbox}{Baby Scale 能支持与不能支持的结论}
它支持：小模型可以从单个儿童环境的转录中学到非随机语言规律；数据来源差异会留下任务相关痕迹。它不支持：儿童输入已经解释人类语言学习；某个家庭 corpus 在大模型上必然最优；相关的数据属性已被证明是因果机制。
\end{importantbox}

评价这种研究还要避免“一个总分覆盖所有能力”。Language-modeling loss 衡量局部预测，BLiMP 类 grammaticality task 检查结构偏好，word similarity 检查 lexical representation，developmental measure 又关心随数据增长的轨迹。某个家庭 corpus 可能在语义任务占优、在句法任务落后；这并非结果矛盾，而是说明数据属性支持不同能力。更强的后续实验应对高相关属性做匹配或干预，例如固定 token 数后重采样说话人多样性、打乱互动 turn、控制词汇覆盖，再看任务轮廓是否按预测变化。这样才能把“哪个数据集更好”推进到“哪种数据机制在起作用”。

样本量与重复训练也必须进入不确定性。小模型实验常受 random seed、batch order 和 early stopping 影响；如果只报告一次 run，数据集排序可能是优化噪声。理想报告应给出多 seed 均值、置信区间、effect size 与 multiple-comparison correction，并说明 family-level transcript 是否存在重叠或质量差异。对于 developmental claim，还应比较数据增长曲线而不是只看最终 checkpoint。这样才能判断某个 corpus 是稳定地更有效，还是刚好在有限训练预算下更快收敛。

\subsection{Bilingual BabyLM：双语是干扰、迁移还是数据规模问题}

双语研究把“混合两种语言”进一步拆成 exposure structure。可以让每个说话人只说一种语言（speaker-split），让说话人在不同句子切换（sentence-level），或在同一句中切换（word-level code-switching）。若把所有条件都叫 bilingual mixture，就无法判断模型到底对语言边界、说话人线索还是总数据量敏感。

\lecturefigure{slide-052.jpg}{Bilingual BabyLM 的三个研究问题：干扰、暴露结构与 scale}{CS25 V6 official deck, p.~52}

\lecturefigure{slide-053.jpg}{训练条件：单语、随机双语、speaker/sentence/word-level switching}{CS25 V6 official deck, p.~53}

设英语与西语数据分布为 $q_{\mathrm{en}}$ 和 $q_{\mathrm{es}}$，混合权重为 $\lambda$，则训练分布为
\[
q_{\lambda}(x)=\lambda q_{\mathrm{en}}(x)+(1-\lambda)q_{\mathrm{es}}(x).
\]
但 speaker-split 与 code-switching 不只是改变 $\lambda$；它们改变语言与上下文、说话人、句界的联合分布。即使总 token 相同，模型可利用的边界线索也不同。

\subsection{结果一：没有普遍双语干扰}

前面的实验设计把双语 exposure 拆成多个条件，本节先回答最基础的问题：同时学习两种语言是否会系统性破坏单语能力。读 slide 55--57 时应先固定总数据与 baseline，再分别检查 language modeling、syntax 和 semantics；只有多个 metric 同方向显著退化，才足以支持广泛 interference。

slide 55 的 PPL 表显示多语模型在两种语言上都能形成有效分布；slide 56--57 再用 grammaticality 与 word similarity 检查是否只是记住局部统计。读表时先找单语 baseline，再看双语条件置信区间，而不是仅比较一位小数。

\lecturefigure{slide-055.jpg}{双语模型在英语与西语上的 perplexity：能同时建模两种语言}{CS25 V6 official deck, p.~55}

\lecturefigure{slide-056.jpg}{Grammaticality：英语 baseline 与多语条件的结构能力对比}{CS25 V6 official deck, p.~56}

\lecturefigure{slide-057.jpg}{Word similarity：语义表示在单语与多语训练中的对比}{CS25 V6 official deck, p.~57}

\begin{knowledgebox}{读图：如何判断 interference}
Interference 不应只定义为某个条件略低于最高分。需要同时看单语 baseline、数据总量、置信区间和多项任务。如果双语模型在英语上略降却学会西语，总体 tradeoff 与“灾难性干扰”不同；如果固定总数据后仍显著下降，才更接近竞争性容量或优化问题。
\end{knowledgebox}

\subsection{结果二：暴露结构影响弱于预期}

不同 multilingual conditions 的 grammaticality 与 word similarity 曲线整体接近，说明模型并不强依赖“一个说话人一种语言”这种人类常见线索。这个结果应解释为：在当前模型、数据量和 metric 下，exposure structure 的边际影响有限；不能推出 code-switching 在所有语言、规模或生成任务中都无关。

\lecturefigure{slide-059.jpg}{不同双语暴露结构下的 grammaticality 表现}{CS25 V6 official deck, p.~59}

\lecturefigure{slide-060.jpg}{不同双语暴露结构下的 word similarity 表现}{CS25 V6 official deck, p.~60}

\subsection{结果三：总数据规模比小幅模型扩张更关键}

slide 62 将约二十分之一参数规模与约二十分之一数据规模分开比较。结果强调 data scale 的作用：减少数据会显著恶化多项指标，而在此小模型范围内， modest model-size change 的影响更小。这里仍需避免外推到 frontier LLM；在更大规模下，参数、数据和 compute 的最优比例可能不同。

\lecturefigure{slide-062.jpg}{数据规模与模型规模的对照：减少数据造成更明显退化}{CS25 V6 official deck, p.~62}

\lecturefigure{slide-063.jpg}{RQ1 结论：没有发现普遍的双语干扰}{CS25 V6 official deck, p.~63}

\lecturefigure{slide-064.jpg}{RQ2 结论：speaker 与 code-switching 结构影响小于预期}{CS25 V6 official deck, p.~64}

\lecturefigure{slide-065.jpg}{RQ3 结论：当前范围内 data scale 比 modest model size 更重要}{CS25 V6 official deck, p.~65}

\lecturefigure{slide-066.jpg}{Bilingual BabyLM takeaways 与实验边界}{CS25 V6 official deck, p.~66}

\teachervoice{00:25:25--00:29:01，老师总结三点：没有广泛 monolingual interference；暴露结构的影响出乎意料地弱；总数据规模比当前范围内的模型大小更关键。老师同时把结论限定在 small-scale setting。}

\begin{warningbox}{“双语不干扰”不是无条件定律}
语言距离、tokenizer、混合比例、训练时长、生成任务与模型容量都会改变结果。当前实验显示的是一组受控 small-scale 条件下没有系统性退化，而不是所有 multilingual training 都自动正迁移。
\end{warningbox}

对工程实践而言，这组结果更像一个调试顺序，而不是 mixture recipe。若 multilingual model 退化，先检查每种语言的有效 token 数、tokenizer fertility、batch sampling 与 evaluation leakage，再检查 capacity competition。Speaker-split 或 code-switching structure 可能在更强语境任务中变得重要；当前语法和词相似指标对 discourse-level language switching 并不敏感。课程的谨慎之处在于把“没有观察到显著 effect”与“effect 不存在”分开，这也是阅读所有 negative result 的基本原则。

\subsection{本章小结}

两个 BabyLM 项目共同说明：small-scale 实验的价值在于分离数据变量。Baby Scale 显示家庭数据属性与任务表现相关；Bilingual BabyLM 显示语言混合不必然造成干扰，且暴露结构与总数据量的作用不同。它们提供机制线索，而不是把儿童学习直接复制到大模型的万能配方。

